%-------------------------
% Illya Yurchanka - Improved CV
%------------------------

\documentclass[11pt]{article}

\usepackage[T1]{fontenc}
\usepackage{helvet}
\renewcommand*\familydefault{\sfdefault}

\usepackage{geometry}
\geometry{
a4paper,
top=2.0cm,
bottom=0.5in,
left=2.5cm,
right=2.5cm
}

\setcounter{secnumdepth}{0}
\pdfgentounicode=1

\usepackage{enumitem}
\setlist[itemize]{
    noitemsep,
    left=0pt..1.5em
}
\setlist[description]{itemsep=0pt}
\setlist[enumerate]{align=left}

\usepackage[dvipsnames]{xcolor}
\colorlet{icnclr}{gray}

\usepackage{titlesec}
\titlespacing{\subsection}{0pt}{*0}{*0}
\titlespacing{\subsubsection}{0pt}{*0}{*0}
\titleformat{\section}{\color{Sepia}\large\bfseries\uppercase}{}{}{\ruleafter}[\global\RemVStrue]
\titleformat{\subsection}{\large\bfseries}{}{}{\rvs}
\titleformat{\subsubsection}{\normalsize}{}{}{}

\usepackage{xhfill}
\newcommand\ruleafter[1]{#1~\xrfill[.5ex]{1pt}[gray]}

\newif\ifRemVS
\newcommand{\rvs}{
    \ifRemVS
        \vspace{-1.5ex}
    \fi
    \global\RemVSfalse
}

\usepackage[bookmarks=false]{hyperref}
\hypersetup{
    colorlinks=true,
    urlcolor=Sepia,
    pdftitle={Illya Yurchanka - CV},
}

\usepackage{fancyhdr}
\pagestyle{fancy}
\renewcommand{\headrulewidth}{0pt}
\fancyhf{}


\begin{document}

%== HEADER ==%
\begin{center}
    {\fontsize{25}{25}\selectfont\textbf{Illya Yurchanka}} \\ \bigskip
    \href{mailto:i.yurchanka@gmail.com}{i.yurchanka@gmail.com} $|$
    \href{https://github.com/illyayurchanka}{github.com/illyayurchanka} $|$
    +48 572 473 066 $|$ Warsaw, Poland
\end{center}

\vspace{0.5ex}
\begin{center}
\parbox{0.92\textwidth}%
\small Physics M.Sc. student at the University of Warsaw with strong programming skills in Java, C++, Python, and SQL. Experience building production-ready software systems and applying data-driven analysis to solve complex scientific problems. Passionate about writing clean, efficient code and delivering end-to-end solutions that empower organizations to make informed decisions. Seeking software engineering roles where I can contribute to mission-critical data platforms.
\end{center}
\vspace{0.5ex}

\section{Skills}
\begin{description}
    \item[\textbf{Programming:}] Python, Java, C++ (STL, Eigen, Boost), SQL (MySQL, PostgreSQL, DuckDB), JavaScript/TypeScript
    \item[\textbf{Data Structures \& Algorithms:}] Trees, Graphs, Hash Tables, Sorting, Searching, Dynamic Programming, Complexity Analysis
    \item[\textbf{Data Platforms \& Engineering:}] Spark, Pandas, NumPy, Scikit-Learn, PyTorch, ML Pipelines, ETL, Data Processing
    \item[\textbf{Tools \& DevOps:}] Git, Docker, Linux/Bash, CMake, Make, ROOT/TMVA, PyTest, CI/CD
    \item[\textbf{Languages:}] English C1, Polish C1, Russian (Native), Belarusian (Native)
\end{description}

\section{Education}
\subsection{University of Warsaw $|$ {\normalfont\textit{M.Sc. in Physics}} \hfill Oct 2025 -- Jun 2027 (expected)}
\begin{itemize}
    \item GPA: 4.5/5 \quad Focus: Particle Physics \& Machine Learning applications
\end{itemize}
\subsection{University of Warsaw $|$ {\normalfont\textit{B.Sc. in Physics}} \hfill Oct 2021 -- Jun 2025}
\begin{itemize}
    \item GPA: 4.5/5 \quad Thesis: Identification of $\Lambda$ hyperons in CBM spectrometer using TMVA (ML)
\end{itemize}

\section{Experience}
\subsection{M.Sc. Researcher -- CBM Experiment, GSI/FAIR (Darmstadt) \hfill Oct 2025 -- Present}
\begin{itemize}
    \item Built and owned full data processing pipeline from ingestion to analysis output; designed and implemented C++ workflows for the CBM detector used in high-energy physics experiments.
    \item Applied ML techniques to simulated Au+Au collision data to improve particle identification and background rejection; results inform detector commissioning strategy.
    \item Collaborated with international team of 15+ researchers to integrate analysis modules into the experiment's core software framework.
\end{itemize}

\subsection{Research Intern -- General Relativity \hfill Mar 2025 -- Jun 2025}
\subsubsection{Center for Theoretical Physics, Polish Academy of Sciences}
\begin{itemize}
    \item Conducted independent research on general relativity; reviewed 10+ peer-reviewed publications to map the state-of-the-art and identify gaps for novel theoretical contributions.
    \item Built mathematical models and performed symbolic/numerical computations in Mathematica and SciPy, producing analytical results presented to senior researchers.
\end{itemize}

\subsection{Goldman Sachs Warsaw Hackathon -- Quant Challenge \hfill Dec 2024}
\begin{itemize}
    \item Led end-to-end implementation of Monte Carlo simulation framework for financial market forecasting; delivered complete solution from design to final presentation.
    \item Implemented variance-reduction techniques -- bootstrapping and importance sampling -- achieving measurably tighter confidence intervals than competing teams.
    \item Collaborated with team members to integrate quantitative models into unified analysis pipeline under tight deadline.
\end{itemize}

\subsection{Warsaw Econometric Challenge \hfill Apr 2025}
\begin{itemize}
    \item Led a team of 4 in a data science competition; rapidly learned SQL and DuckDB to deliver ML-based predictive solutions under strict deadline.
    \item Built end-to-end ML pipeline using Scikit-Learn -- from data ingestion and feature engineering to model training and predictions.
    \item Presented results to judges and answered technical questions about methodology and implementation choices.
\end{itemize}

\section{Projects}

\subsection{Multi-Head Latent Attention Implementation \hfill Oct 2025}
\textit{Python, PyTorch} \quad \href{https://github.com/illyayurchanka/nanodeepsk}{\small[GitHub]}
\begin{itemize}
    \item Implemented MLA mechanism from scratch in PyTorch; engineered efficient KV cache compression through low-dimensional latent space projection.
    \item Optimized for memory efficiency while preserving attention quality; benchmarked against standard multi-head attention.
\end{itemize}

\subsection{Food Nutrition Analysis with Machine Learning \hfill Jan 2026}
\textit{Python, Scikit-Learn, Pandas, NumPy, Transformers} \quad \href{https://github.com/illyayurchanka/nutritions}{\small[GitHub]}
\begin{itemize}
    \item Built production-ready ML pipeline on the Open Food Facts dataset: multi-task learning for NutriScore prediction, calorie estimation, NOVA classification, and Eco-Score analysis.
    \item Implemented semantic search engine using vector embeddings for product recommendation; optimized for retrieval speed.
\end{itemize}

\subsection{Bachelor Thesis -- $\Lambda$ Hyperon Identification in CBM Spectrometer \hfill May 2025}
\textit{C++, ROOT, ROOT-TMVA}
\begin{itemize}
    \item Benchmarked KNN, AdaBoost, and MLP classifiers for signal/background separation in high-energy physics data; ensemble methods achieved the best ROC-AUC, directly informing the experiment's offline reconstruction strategy.
\end{itemize}

\subsection{Terminal Cookie Clicker Game \hfill Apr 2024}
\textit{C++, Ncurses}
\begin{itemize}
    \item Built a real-time terminal incremental game with live input handling, persistent game-state management, and a full ncurses TUI rendered from scratch.
\end{itemize}

\section{Relevant Courses}
\begin{description}
    \item[\textbf{Machine Learning}] Neural networks (FC, CNN, GAN, VAE), RNNs, Transformers, NumPy/Pandas/Matplotlib \hfill 2023/2024
    \item[\textbf{Programming \& Numerical Methods}] Advanced C++, physical simulations, numerical methods in Python \hfill 2023/2024
    \item[\textbf{Statistical Analysis of Experimental Data}] Statistical inference, hypothesis testing, data analysis pipelines \hfill 2025/2026
    \item[\textbf{Programming}] STL algorithms, templates, containers \hfill 2022/2023
\end{description}

\end{document}
